\subsection{可数集}

定义3. 如果一个集合与整数集N的某个子集等势，则称该集合为可数的。

命题1. 可数集的每个子集都是可数的。有限个可数集的笛卡尔积是可数的。一列可数集的并集是可数的。

第一个断言是显然的。其余断言由第3号中定理2的各推论得出。

我们已证明（第3号，引理1）若 \(\mathfrak{a}\) 是任意无限基数，则 \(\operatorname{Card}(\mathbf{N}) \leqslant \mathfrak{a}\) 。这有以下推论：

命题2. 每个可数无限集 E 与 N 等势。

根据定义有 \(\operatorname{Card}(\mathbf{E}) \leqslant \operatorname{Card}(\mathbf{N})\) ；且有 \(\operatorname{Card}(\mathbf{N}) \leqslant \operatorname{Card}(\mathbf{E})\) ，因为 \(\mathbf{E}\) 是无限集。

命题3. 每个无限集都有一个划分 \((\mathbf{X}_i)_{i\in I}\) ——由可数无限集 \(\mathbf{X}_i\) 组成——其指标集 I 与 E 等势。

因为 \(\operatorname{Card}(\mathbf{E}) = \operatorname{Card}(\mathbf{E})\) \(\operatorname{Card}(\mathbf{N})\) （第 3 号，定理 2，推论 4）。

命题4。设 \(f\) 是集合 \(\mathbf{E}\) 到无穷集 \(\mathbf{F}\) 上的一个映射，使得对于每个 \(y \in \mathbf{F}\) , \(\overline{f}^1(y)\) 是可数的。那么 \(\mathbf{F}\) 与 \(\mathbf{E}\) 等势.

因为集合 \(\overline{f}^1 (y)\) \((y\in \mathbf{F})\) 构成 \(\mathbf{E}\) 的一个划分；因此

\[
\operatorname{Card} (\mathbf {E}) \leqslant \operatorname{Card} (\mathbf {F}) \operatorname{Card} (\mathbf {N}) = \operatorname{Card} (\mathbf {F}),
\]

并且 \(\operatorname{Card}(\mathbf{F}) \leqslant \operatorname{Card}(\mathbf{E})\) 由第3节第2号的命题3，

命题5. 集合 \(\mathfrak{F}(\mathbf{E})\) ——无限集 \(\mathbf{E}\) 的有限子集所成的集合——与 \(\mathbf{E}\) 等势.

对每个整数 \(n\) ，让 \(\mathfrak{F}_n\) 表示 \(\mathbf{E}\) 的具有 \(n\) 个元素的子集所组成的集合。对每个 \(\mathbf{X} \in \mathfrak{F}_n\) ，存在 \([1, n]\) 到 \(\mathbf{X}\) 上的一个双射。因此， \(\mathfrak{F}_n\) 的基数至多等于 \([1, n]\) 到 \(\mathbf{E}\) 中的映射所成集合的基数，即 \(\operatorname{Card}(\mathbf{E}^n) = \operatorname{Card}(\mathbf{E})\) (第3号，定理2，推论1)。因此

\[
\operatorname{Card} (\mathfrak {F} (E)) = \sum_ {n \in N} \operatorname{Card} (\mathfrak {F} _ {n}) \leqslant \operatorname{Card} (E) \operatorname{Card} (N) = \operatorname{Card} (E).
\]

另一方面，由于 \(x \to \{x\}\) 是 \(\mathbf{E}\) 到 \(\mathfrak{F}(\mathbf{E})\) 中的一个单射映射，我们有 \(\operatorname{Card}(\mathbf{E}) \leqslant \operatorname{Card}(\mathfrak{F}(\mathbf{E}))\) .

定义4. 一个集合被称为具有连续统的势，如果它与N的所有子集的集合等势。

一个具有连续统势的集合是不可数的（§ 3，第 6 号，定理 2）。

* “连续统的势”这一名称源于实数集与 \(\mathfrak{P}(\mathbf{N})\) 等势这一事实（一般拓扑学，第四章，第8节）。* 连续统假设是指每个不可数集合都包含一个具有连续统势的子集；广义连续统假设是指，对于每个无限基数 \(a\) ，每个基数 \(>a\) 是 \(\geqslant 2^{a}\) .

